\chapter{Sensado humano mediante Wi-Fi y LoRa}
\label{ch:sensado_wifi_lora}

El sensado inalámbrico reutiliza infraestructura de comunicación como instrumento
indirecto de la escena. La literatura demuestra viabilidad en presencia, actividad
y respiración, pero también muestra que el desempeño depende del observable, la
geometría y el dominio de evaluación. Este capítulo organiza esa evidencia según
la variable física y el mecanismo de información que la hace recuperable.

\section{Presencia, movimiento y actividad}

La presencia es una decisión entre distribuciones del canal con y sin persona; la
actividad añade una trayectoria temporal del cuerpo. Ambas tareas suelen producir
cambios mayores que la respiración, pero involucran oclusiones, variación de soporte
y múltiples segmentos. Por ello, un clasificador puede explotar regularidades de
escena sin haber aprendido un mecanismo transferible. Las revisiones recientes
identifican la variación entre habitaciones, sujetos, dispositivos y protocolos
como una limitación central \citep{Zhu2024Survey,Armenta2024Survey}.

La magnitud, diferencias de fase, Doppler y correlaciones espaciales ofrecen vistas
complementarias. La magnitud es menos sensible a una referencia común, pero puede
perder pequeños cambios rotacionales; la fase conserva micromovimiento si el reloj
es estable; las ventanas Doppler separan escalas temporales a costa de asumir
estacionariedad local. La elección debe justificarse por el estado $q$ y no sólo
por exactitud de clasificación.

Wi-Actigraph estudia actividad nocturna y eventos fisiológicos en un entorno
realista, mientras el trabajo de Galdino y colaboradores sitúa Wi-Fi sensing dentro
de salud ubicua \citep{Alzaabi2025WiActigraph,Galdino2023eHealth}. Estos antecedentes
justifican infraestructura no invasiva. Sus etiquetas y protocolos no acreditan
por sí mismos una forma de onda fisiológica ni transferencia a otra cadena de
adquisición; establecen tareas downstream sobre observables específicos.

\section{Respiración y signos vitales como problema físico}

En primera aproximación, la respiración induce un desplazamiento torácico
$x_{\rm chest}(t)$ que perturba longitudes y coeficientes de varias trayectorias.
El objetivo más informativo es recuperar una forma de onda alineada con una
referencia. La frecuencia respiratoria es una cantidad derivada y puede ocultar
errores de amplitud, fase temporal, apneas o movimiento. Sólo después de validar
esa cadena procede estudiar variables clínicas.

BreatheSmart combina un maniquí respiratorio y experimentos humanos, y muestra
que amplitud y fase de CSI contienen información dependiente de posición y
configuración \citep{Mosleh2022BreatheSmart}. Su valor para esta tesis reside en
el uso de una perturbación controlada y en documentar fallas geométricas, no en
trasladar directamente su desempeño. FarSense forma una razón entre antenas para
rechazar perturbaciones comunes y extender alcance \citep{Zeng2019FarSense}. La
razón requiere un denominador con energía suficiente y puede eliminar un cambio
físico común; será tratada como una hipótesis de invariancia que debe contrastarse.

El movimiento concurrente no es ruido aditivo independiente. Puede cambiar el
soporte de trayectorias y modular con mayor energía la misma medición que contiene
respiración. MoVi-Fi, aunque utiliza radar IR-UWB/FMCW, ilustra la separación bajo
movimiento \citep{Chen2022MoViFi}; ViMo aporta evidencia a \SI{60}{GHz}
\citep{Wang2020ViMo}. Sus ventajas de ancho de banda y apertura no son transferibles
sin más a CSI sub-7 GHz. La comparación debe declarar frecuencia, $B$, apertura,
cadencia, referencia temporal y libertad de movimiento.

\section{Múltiples usuarios como separabilidad física}

Con dos personas, el canal depende de $\vect q=[\vect q_1^T,\vect q_2^T]^T$. Aun
si cada sujeto es observable de manera aislada, los Jacobianos pueden ser casi
colineales:
\begin{equation}
\mat J_q=[\mat J_{q_1}\ \mat J_{q_2}].
\end{equation}
La separación requiere que retardo, ángulo, enlace o evolución temporal aporten
direcciones distintas. Dos tasas respiratorias bien separadas ayudan a un método
espectral, pero no constituyen una condición general; frecuencias cercanas,
armónicos, no estacionariedad y asociación persona--componente generan ambigüedad.

WiMANS combina CSI y video sincronizado para actividades de cero a cinco personas
y documenta el deterioro al aumentar ocupación \citep{Huang2024WiMANS}. WiMUSE
formula respiración multiusuario como separación ciega y expone límites de métodos
puramente espectrales \citep{Yi2025WiMUSE}. MURAL-Fi incorpora diversidad de AoA,
ToF, AoD y Doppler para separar regiones torácicas y tolerar movimiento involuntario
\citep{Li2026MURALFI}. La lección conjunta es física: el aprendizaje organiza la
diversidad disponible, pero no crea una dimensión ausente.

UWB-Fi amplía el ancho de banda efectivo mediante muestreo disperso en frecuencia
y combina recuperación estructurada con conocimiento de trazado de rayos
\citep{Li2024UWBFI}. Esta estrategia mejora resolución potencial a cambio de
coherencia entre canales, tiempo de barrido y una cadena distinta de Wi-Fi
convencional. HiSAC estudia la misma motivación general desde señales de
comunicación multibanda y sirve como referencia para distinguir ancho de banda
agregado de ancho de banda coherentemente utilizable \citep{Pegoraro2024HiSAC}.
Ninguna de estas ventajas se atribuirá a una CSI monobanda sin reproducir su
operador de adquisición. Por ello, la tesis evalúa primero una fuente y habilita el régimen
multiusuario sólo si la información conjunta predice separabilidad y permite
abstención en geometrías ambiguas.

\section{Generalización y tipos de transferencia}

La palabra transferencia reúne problemas diferentes. La transferencia espacial
reserva posiciones; la temporal reserva sesiones; la instrumental cambia radio o
firmware; la ambiental cambia escena; y la poblacional reserva sujetos. Una
partición aleatoria de ventanas correlacionadas no prueba ninguno de estos casos.
Las unidades completas del dominio objetivo deben quedar fuera del ajuste.

MM-Fi proporciona radio y referencias multimodales sincronizadas para estudiar
correspondencias \citep{Yang2023MMFi}. La visión puede aportar supervisión,
geometría o desambiguación durante el desarrollo \citep{Hori2024Fusion}; debe
declararse si está disponible en prueba. Una ganancia que requiere cámara no puede
atribuirse exclusivamente a RF. De modo análogo, una posición exacta usada por el
gemelo es una entrada adicional y debe compararse con alternativas que reciben
información equivalente.

WiFi-JEPA propone aprendizaje predictivo en espacio latente
\citep{Kim2026WiFiJEPA}. El preentrenamiento puede reducir etiquetas si el objetivo
obliga a conservar dinámica útil. También puede codificar identidad de escena o
descartar cambios pequeños. Por ello se evaluará mediante curvas de error frente
a cantidad de adaptación y pruebas de preservación del Jacobiano, además de una
métrica downstream. La física funciona como condición y diagnóstico; no sustituye
la evaluación fuera de dominio.

\section{LoRa: de CSS al observable de sensado}

LoRa usa \ac{CSS}: la frecuencia instantánea recorre un ancho de banda durante cada
símbolo. Para una envolvente $x_{\rm CSS}$,
\begin{equation}
y_{\rm L}(t)=\sum_{\ell}\alpha_{{\rm L},\ell}(t)
x_{\rm CSS}(t-\tau_{{\rm L},\ell}(t))+w_{\rm L}(t).
\label{eq:lora_channel}
\end{equation}
El receptor multiplica por un chirp conjugado y aplica una DFT. Bajo el modelo
ideal, el \emph{dechirp} convierte el símbolo en un tono cuya localización permite
demodulación \citep{Ghanaatian2019LoRaRx,Maleki2023CSS}. El vector complejo antes
de seleccionar el bin contiene más información que RSSI, SNR o el símbolo
decodificado.

CFO, STO y SFO no desaparecen con el dechirp. Desplazan o ensanchan picos, introducen
rotaciones entre símbolos y pueden confundirse con dinámica lenta. La estimación
para comunicación busca decodificar; el sensado requiere además estabilidad o una
referencia para atribuir pequeñas variaciones al entorno. Una pasarela LoRaWAN
convencional suele entregar paquetes y resúmenes, no necesariamente I/Q coherente.
Por ello se distingue acceso de pasarela, recepción SDR y hardware con muestras
intermedias.

Zhang y colaboradores demuestran sensado de largo alcance y a través de paredes
mediante una razón entre antenas \citep{Zhang2020LoRa}. SenLoRa aprovecha símbolos
de alta confianza de tráfico ambiente y convierte la disponibilidad temporal en
una restricción conjunta de comunicación y sensado \citep{SenLoRa2025}. El estudio
de gran escala de Xue y colaboradores relaciona configuración del transceptor con
cobertura efectiva de sensing \citep{Xue2025LargeScaleLoRa}. Estas obras respaldan
alcance y factibilidad, pero no garantizan que una pasarela particular preserve
fase ni que su cadencia baste para respiración.

\subsection{Calidad de comunicación y calidad para sensado}

Para $Y(t)=H_s+H_d(t)+n(t)$, sean
$P_d=\E|H_d|^2$ y $P_n=\E|n|^2$. La SNR empleada para comunicación,
\begin{equation}
\mathrm{SNR}_{\rm comm}
=\frac{\E|H_s+H_d|^2}{P_n},
\end{equation}
responde si la señal total puede detectarse o demodularse. En cambio,
\begin{equation}
\mathrm{SSNR}=\frac{P_d}{P_n},\qquad
\eta_d=\frac{P_d}{\E|H_s+H_d|^2},
\label{eq:ssnr_dynamic_fraction}
\end{equation}
describen la componente dinámica sobre ruido y su fracción respecto del canal.
Un LOS estático intenso puede producir $\mathrm{SNR}_{\rm comm}$ elevada y
$\eta_d$ pequeña; el enlace comunica con fiabilidad y, sin embargo, apenas cambia
con la persona.

Bajo una suma \emph{incoherente} de potencias,
$\eta_d\mathrm{SNR}_{\rm comm}\approx\mathrm{SSNR}$. La aproximación no es general,
pues el término $2\Re\{H_sH_d^*\}$ depende de fase. SenLoRa utiliza además una
relación de contraste de sensing de la forma
\begin{equation}
\mathrm{SCR}=
\frac{\E\left|H_d/(H_s+H_d)\right|^2}{\sigma_n^2},
\label{eq:senlora_scr}
\end{equation}
bajo la normalización de ruido adoptada en su procesamiento \citep{SenLoRa2025}.
La normalización debe declararse para que la escala sea interpretable. El cociente
puede crecer cerca de $H_s+H_d\simeq0$, precisamente donde la fase, el control de
ganancia y la demodulación son menos estables. También requiere estimar una
descomposición $H_s/H_d$ que no se observa directamente. Por estas razones, SSNR,
$\eta_d$ y SCR son métricas intermedias y no pruebas de recuperación respiratoria.

\subsection{Suficiencia temporal y entropía de muestra}

SenLoRa emplea entropía de muestra para evaluar la estabilización de un descriptor
temporal. Para patrones de longitud $m$ y tolerancia $r$,
\begin{equation}
\operatorname{SampEn}(m,r)=
-\log\frac{A_{m+1}(r)}{B_m(r)},
\label{eq:sample_entropy}
\end{equation}
donde $B_m$ cuenta coincidencias de longitud $m$ y $A_{m+1}$ las que permanecen
coincidentes al añadir una muestra. La convergencia reportada alrededor de 600
chirps pertenece a su configuración; no constituye un umbral universal. Depende
de factor de ensanchamiento, ancho de banda, duración y número de chirps, tráfico,
$m$, $r$, ruido y dinámica objetivo.

Además, alta entropía puede provenir de ruido y baja entropía de una señal constante
o de respiración periódica limpia. La tesis utilizará la entropía de muestra como
diagnóstico de regularidad y suficiencia de ventana, distinguiéndola de entropía de
Shannon o información de Fisher. Para forma de onda se reportarán duración
$T_{\rm obs}$, ciclos observados, cadencia efectiva y distribución de huecos; cumplir
Nyquist para la frecuencia fundamental no garantiza resolución de inspiración,
espiración y pausas.

OpenLoRa ofrece un banco para recepción PHY e interferencia
\citep{Mishra2023OpenLoRa}. En esta tesis se utiliza para auditar el operador y
las perturbaciones, no como evidencia de sensado humano. La cadena de conservación
se documentará como
\[
\text{I/Q}\rightarrow\text{I/Q tras dechirp}\rightarrow
\text{DFT compleja}\rightarrow(A,\phi)\rightarrow
(\mathrm{RSSI},\mathrm{SNR}).
\]
Interferencia, colisiones y tráfico esporádico afectan tanto calidad por muestra
como disponibilidad. Un resumen puede regularizar un estimador con datos finitos,
pero una transformación no invertible no incrementa la información del modelo
original \citep{Kay1993,CoverThomas2006}.

\section{Dos radios y una pregunta común}

Wi-Fi y LoRa comparten el estado físico, pero implementan operadores distintos:
\begin{equation}
Y_{\rm WiFi}=\mathcal O_{\rm WiFi}(H_{\rm WiFi}),\qquad
Y_{\rm LoRa}=\mathcal O_{\rm LoRa}(H_{\rm LoRa}).
\end{equation}
Wi-Fi ofrece inicialmente mayor diversidad multicarrier para microdinámica; LoRa
ofrece presupuesto de enlace, cobertura y otra relación entre tráfico y muestreo.
La comparación pertinente no pregunta cuál tecnología es universalmente superior,
sino cuánta sensibilidad útil preserva cada operador bajo igual perturbación y
qué costo instrumental requiere.

La literatura deja tres brechas conectadas: separar robustez instrumental de
sensibilidad física, evaluar un gemelo por su respuesta y no sólo por su campo
estático, y declarar cuándo la recuperación deja de ser identificable. El capítulo
siguiente formaliza estas brechas mediante invariancia, resolubilidad, Jacobianos
e información de Fisher.
